%----------------Cap_04----------------%
\chapter{Conclusão}
 
Este trabalho propôs e documentou o desenvolvimento de um controlador neural
embarcado que aproxima a política ótima de um Controle Preditivo Não-Linear
(NMPC) para rastreamento angular de um aeropêndulo instrumentado. A motivação
central reside na incompatibilidade entre o custo computacional da otimização
repetida do NMPC e os recursos disponíveis em microcontroladores de baixo custo,
problema recorrente em plataformas VTOL de pequeno porte.

Como referência de desempenho e como mecanismo de estabilização durante os
ensaios de identificação, um controlador PID foi sintonizado a partir do modelo
linearizado e implementado digitalmente com discretização trapezoidal e
anti-windup por \textit{back-calculation}. Os resultados preliminares da
Etapa~\ref{subsec:etapa2} confirmaram a não-linearidade do sistema, motivando o emprego do modelo não-linear para a síntese do NMPC.
 
A estratégia adotada é tecnicamente viável. Os três principais riscos identificados, divergência do modelo identificado, exaustão dos recursos de hardware e
instabilidade do controlador neural por dados fora da distribuição, possuem
ações preventivas e de contingência estabelecidas. A etapa de identificação (Etapa~\ref{subsec:etapa3})
encontra-se em andamento e as etapas subsequentes dependem diretamente da
qualidade do modelo obtido. Caso o critério de ajuste não seja satisfeito, a
estrutura do modelo será revisada antes de prosseguir. A compacidade da arquitetura
MLP e a quantização via X-CUBE-AI constituem caminhos de contingência caso
os limites de Flash e SRAM do STM32F446 sejam ultrapassados.
 
O cronograma prevê a conclusão da identificação e da síntese do NMPC no
primeiro semestre, reservando o segundo para o treinamento da rede, a implantação
embarcada e os ensaios experimentais comparativos. A comparação final entre PID,
NMPC e MLP embarcada, consolidada na Tabela~\ref{tab:comparacao_controladores},
constituirá a principal contribuição quantitativa do trabalho e permitirá avaliar
objetivamente os \textit{trade-offs} entre desempenho de rastreamento, rejeição a
perturbações e custo computacional em hardware embarcado de baixo custo.
 